\documentclass[11pt]{article}

\usepackage[utf8]{inputenc}
\usepackage[T1]{fontenc}
\usepackage{lmodern}
\usepackage{amsmath}
\usepackage{amssymb}
\usepackage{amsthm}
\usepackage[hidelinks]{hyperref}
\usepackage{doi}
\hypersetup{
  unicode,
  pdftitle={Projected Yukawa Line and Generation-Level Assignment: Functorial Closure of the Determinant Line
            and Spectral Ordering on the Generation Triplet},
  pdfauthor={Jérôme Beau},
  pdfkeywords={projected Yukawa sector; determinant line; admissible spinor bundle; tensor functors; hypercharge line;
               generation ordering; exit deficit; squared projective residue; spectral level assignment;
               non-injective projection}
}

\input{glyphtounicode}
\pdfgentounicode=1

\newtheorem{theorem}{Theorem}
\newtheorem{proposition}{Proposition}
\newtheorem{lemma}{Lemma}
\newtheorem{corollary}{Corollary}
\theoremstyle{definition}
\newtheorem{definition}{Definition}
\theoremstyle{remark}
\newtheorem{remark}{Remark}

\newcommand{\JPi}{J_{\Pi}}
\newcommand{\Jthree}{J_{3}}
\newcommand{\Rmix}{R_{\mathrm{mix}}}
\newcommand{\Cgen}{\mathbb{C}^{3}_{\mathrm{gen}}}
\newcommand{\Epi}{E_{\Pi}}
\newcommand{\Sbundle}{\mathcal{S}_{\Pi}}
\newcommand{\Vr}{\mathcal{V}_{\mathbb{R}}}
\newcommand{\LY}{L_{Y}}
\newcommand{\Sym}{\mathrm{Sym}}
\newcommand{\adC}{\mathrm{ad}_{\mathbb{C}}}
\newcommand{\SU}{\mathrm{SU}}
\newcommand{\UU}{\mathrm{U}}
\newcommand{\SL}{\mathrm{SL}}
\newcommand{\GL}{\mathrm{GL}}
\newcommand{\Spin}{\mathrm{Spin}}
\newcommand{\HS}{\mathrm{HS}}
\newcommand{\tr}{\operatorname{tr}}
\newcommand{\diag}{\operatorname{diag}}

\title{
  Projected Yukawa Line and Generation-Level Assignment
  \\[1ex]
  \large Functorial Closure of $\LY = \wedge^{2}(\Sbundle)$ and Spectral Ordering on $\Cgen$
}

\author{J. Beau\\
Independent Researcher, France\\
\texttt{jerome.beau@cosmochrony.org}}

\date{2026}

\begin{document}

  \maketitle

  \begin{abstract}
    This note opens the mass-sector frontier of the fermionic-matter sub-programme of Cosmochrony, and closes only its
    first stage.
    It does not derive the three-generation masses, the Yukawa couplings, or the mixing.
    It isolates two structural results that the admissible spinor carrier already determines, and separates them sharply
    from the mass normalisation that does not yet follow.
    First, on the rank-two admissible spinor bundle $\Sbundle$ produced by the Lorentzian complexification of the
    metaplectic Weil carrier of Q14, the determinant line $\LY := \wedge^{2}(\Sbundle)$ is the unique
    functorial line of $\Sbundle$ that is invisible to the adjoint sector $\Sym^{2}(\Sbundle) \simeq \adC(P_{\Pi})$, and
    it is the carrier of the abelian electroweak weight; the projected Yukawa sector is therefore the determinant-line
    sector of the admissible spinor functorial closure, not an external bundle datum.
    Second, on the gauge-singlet generation triplet $\Cgen = \mathrm{span}(e_{0}, e_{+}, e_{-})$ the squared projective
    residue takes the closed normal form $\Epi^{2}|_{\Cgen} = \diag(1, \tfrac12 + u, \tfrac12 - u)$ of the Schur
    reduction (PRS), whose exit deficits $\{0, \tfrac12 - u, \tfrac12 + u\}$ order the three generation levels (Q14).
    The single split coefficient $u$ fixes the dimensionless level structure but not the absolute scale: the projected
    Yukawa operator, its norm, the sign of $u$, and the complex mixing phase remain downstream open data, so no mass
    value is claimed.
    All representation-theoretic and spectral identities are verified by exact symbolic computation, with no sampling.
  \end{abstract}

  \noindent{\small\textbf{Keywords:}
  projected Yukawa sector; determinant line; admissible spinor bundle; tensor functors; hypercharge line; generation
  ordering; exit deficit; squared projective residue; spectral level assignment; non-injective projection.}

  \newpage

  \tableofcontents

  \newpage

  \section{Position and purpose}
    \label{sec:position}

    The fermionic-matter sub-programme identifies the electroweak and generation structure of matter from the admissible
    spinor bundle $\Sbundle$ alone, without an external matter bundle.
    Three identifications are already established in the constituent paper~\cite{Beau2026q14}: the tensor functors of
    $\Sbundle$ reproduce the $\SU(2)_{L} \times \UU(1)_{Y}$ electroweak bundle (Theorem~A), the projected Dirac residue
    $\Epi$ enforces the $V-A$ chirality and fixes hypercharge as a spectral coherence datum (Theorem~B), and the
    saturation invariant forces a gauge-singlet three-generation factor $\Cgen$ (Theorem~C).
    The split coefficient of $\Cgen$ has since been reduced to a single real number $u$, with its existence $u \neq 0$
    closed in the projective-residue Schur reduction~\cite{Beau2026prs} and the radial factor of its modulus closed in the
    Born--Infeld margin note~\cite{Beau2026bim}.
    The sub-programme note states plainly that the mass spectrum and the Yukawa couplings are not yet derived, and lists
    them among the open deliverables~\cite{Beau2026fmnote}.

    This note takes up the first stage of that open deliverable, and only the first stage.
    It does not compute a mass.
    It closes two structural statements that the carrier already determines --- the functorial Yukawa line and the
    generation-level ordering --- and it isolates, as the genuinely open downstream object, the projected Yukawa operator
    whose norm and phase carry the absolute mass scale and the mixing.
    The method lock is stated once and held throughout: the determinant line $\LY$ fixes the coupling line, the squared
    residue $\Epi^{2}|_{\Cgen}$ fixes the generation levels, and the mass comes after.

  \section{The determinant line $\LY$}
    \label{sec:line}

    The carrier is the rank-two admissible spinor bundle $\Sbundle$, the fundamental bundle of
    $\Spin(3,1) \simeq \SL(2,\mathbb{C})$ obtained from the metaplectic lift of the Weil module $\Vr$ under the Lorentzian
    complexification $\mathrm{mp}(2,\mathbb{R})_{\mathbb{C}} \simeq \mathfrak{sl}_{2}(\mathbb{C})$ forced by
    $g^{\mu\nu} = 2\eta^{\mu\nu}$~\cite{Beau2026q14}.
    The tensor functors of $\Sbundle$ are taken on the rank-two fibre, not on the doubled Dirac space
    $\mathcal{S}_{L} \oplus \mathcal{S}_{R}$, so that no cross-chiral components are manufactured.

    \begin{proposition}[Projected Yukawa line]
      \label{prop:yukawa-line}
      Let $\Sbundle$ be the rank-two admissible spinor bundle obtained from the Lorentzian complexification of the
      metaplectic Weil carrier.
      Then the determinant line
      \[
        \LY := \wedge^{2}(\Sbundle)
      \]
      is the unique functorial line produced by $\Sbundle$ that is invisible to the $\SU(2)_{L}$ adjoint sector
      $\Sym^{2}(\Sbundle) \simeq \adC(P_{\Pi})$, and it carries the abelian electroweak weight.
      Consequently the projected Yukawa sector is not an additional external bundle datum, but the determinant-line sector
      of the admissible spinor functorial closure.
    \end{proposition}

    \begin{proof}
      On the rank-two fibre $\mathbb{C}^{2}$ the polynomial functors are the symmetric and exterior powers $\Sym^{k}$ and
      $\wedge^{k}$.
      Their dimensions are $\dim \Sym^{2}(\mathbb{C}^{2}) = 3$ and $\dim \wedge^{2}(\mathbb{C}^{2}) = 1$, and
      $\wedge^{2}(\mathbb{C}^{2})$ is the only nontrivial one-dimensional instance, so $\LY = \wedge^{2}(\Sbundle)$ is the
      unique functorial line of $\Sbundle$.
      The symmetric sector matches the complex adjoint, $\Sym^{2}(\mathbb{C}^{2}) \simeq \mathfrak{sl}_{2}(\mathbb{C})$ of
      dimension $3$, whose compact real form is the $\SU(2)_{L}$ gauge sector~\cite{Beau2026q14}.
      For any $g \in \SL(2,\mathbb{C})$, and in particular for the compact real form $\SU(2)_{L} \subset \SL(2,\mathbb{C})$
      selected from $\Sym^{2}(\Sbundle)$, the induced action on the line $\wedge^{2}(\mathbb{C}^{2})$ is multiplication by
      $\det(g) = 1$.
      Hence $\LY$ carries no $\SU(2)_{L}$ charge: it is invisible to the adjoint sector.
      The induced action of a diagonal abelian generator $\diag(y_{1}, y_{2})$ on $\wedge^{2}(\mathbb{C}^{2})$ is
      multiplication by the trace $y_{1} + y_{2}$, so the determinant line is the carrier of the abelian weight, with
      $\UU(1)_{Y}$ realised as the structure group $\GL(1)$ of the line.
      The identification $\wedge^{2}(\Sbundle) = \LY$ is Theorem~A of~\cite{Beau2026q14}, and no electroweak matter bundle
      is introduced beyond the functors of $\Sbundle$.
    \end{proof}

    The determinant line is thus intrinsic to the carrier.
    It is the abelian half of the same functorial closure whose symmetric half is the non-abelian gauge sector, and it
    requires no separate hypercharge bundle.

  \section{Separation from the Yukawa operator}
    \label{sec:separation}

    Proposition~\ref{prop:yukawa-line} fixes a line, not a mass matrix.
    The line $\LY$ is the weight datum of the Yukawa coupling; the mass matrix is a distinct object, a projected morphism
    of the admissible spinor sectors of the schematic form
    \begin{equation}
      \label{eq:yukawa-operator}
      Y_{\Pi}:\ \mathcal{S}_{L,\Pi} \otimes \LY \longrightarrow \mathcal{S}_{R,\Pi},
    \end{equation}
    or equivalently its spectral avatar after compression by the projective residue $\Epi$.
    The line fixes which abelian weight the coupling carries; the operator $Y_{\Pi}$ fixes how strongly, and in which
    generation directions, the chiralities are tied.
    This separation is the central guardrail of the note: the closure of the coupling line does not close the operator.
    The norm of $Y_{\Pi}$, the orientation (sign) of the split coefficient $u$, and the complex metaplectic mixing phase
    are downstream data, recorded as open in the sub-programme inventory~\cite{Beau2026fmnote}; none of them is supplied
    here.

  \section{Generation-level assignment}
    \label{sec:levels}

    The internal structure of the gauge-singlet triplet $\Cgen$ is read from the squared projective residue $\Epi^{2}$,
    the zero-order spectral residue of non-injective projection in the spinorial sector~\cite{Beau2026q14, Beau2026prs}.
    On $\Cgen = \mathrm{span}(e_{0}, e_{+}, e_{-})$ with Cartan generator $\Jthree = \diag(0, 1, -1)$, the Schur reduction
    fixes the normal form $\Epi^{2}|_{\Cgen} = (C_{2} - \Jthree^{2})/C_{2} + u\,\Jthree + v\,\Rmix$ with spin-one Casimir
    $C_{2} = 2$~\cite{Beau2026prs}.
    The CP-even real sector imposes $v = 0$.
    The value $(C_{2} - \Jthree^{2})/C_{2} = \diag(1, \tfrac12, \tfrac12)$ is algebraic at
    $C_{2} = 2$.
    Its reading as the Born--Infeld even sector would require an identification between the
    conditional $3\times3$ model of~\cite{Beau2026a34} on $\Sym^2(V_\rho)$ and
    $(E_\Pi^2)_{\mathrm{even}}$; no such identification is available, so that reading is not used
    here.

    \begin{proposition}[Generation-level assignment]
      \label{prop:generation-level-assignment}
      On the gauge-singlet generation triplet $\Cgen = \mathrm{span}(e_{0}, e_{+}, e_{-})$, the projected Yukawa levels
      are ordered by the spectral levels of
      \[
        \Epi^{2}\big|_{\Cgen}
        =
        \diag\!\left(1,\ \tfrac12 + u,\ \tfrac12 - u\right).
      \]
      With the central $\JPi$-fixed level $s_{0} = 1$, the exit deficits $d_{\mathrm{exit}}(e_{i}) := s_{0} - s_{i}$ are
      \[
        d_{\mathrm{exit}}(e_{0}) = 0,
        \qquad
        d_{\mathrm{exit}}(e_{+}) = \tfrac12 - u,
        \qquad
        d_{\mathrm{exit}}(e_{-}) = \tfrac12 + u,
      \]
      so for $u > 0$ the ordering $0 < d_{\mathrm{exit}}(e_{+}) < d_{\mathrm{exit}}(e_{-})$ assigns $e_{0}$ the lightest and
      $e_{-}$ the heaviest generation level.
      The assignment determines the three generation levels structurally.
      It does not by itself determine the physical mass normalisation, which still depends on the projected Yukawa norm
      and the upstream dictionary data.
    \end{proposition}

    \begin{proof}
      The normal form is~\cite{Beau2026prs} with $v = 0$ and the algebraic even part
      $(C_{2} - \Jthree^{2})/C_{2} = \diag(1, \tfrac12, \tfrac12)$ at $C_{2} = 2$; its eigenvalues on
      $(e_{0}, e_{+}, e_{-})$ are $(1, \tfrac12 + u, \tfrac12 - u)$.
      The exit deficits follow from the support-exit definition relative to the non-exiting central level
      $s_{0} = 1$~\cite{Beau2026q14, Beau2026a4}.
      The ordering for $u > 0$ follows from $d_{\mathrm{exit}}(e_{-}) - d_{\mathrm{exit}}(e_{+}) = 2u > 0$, which places
      the central direction $e_{0}$ at zero deficit (lightest, non-exiting) and $e_{-}$ at the largest deficit (heaviest).
    \end{proof}

    The level structure is dimensionless and fixed by $u$ alone: the ratio $(\tfrac12 + u)/(\tfrac12 - u)$ is invariant
    under the spectral rescaling $\Epi^{2} \mapsto \lambda\,\Epi^{2}$~\cite{Beau2026prs}, so the generation ordering does
    not depend on the absolute spectral normalisation.
    The numerical value of $u$ is itself dictionary-bound through the chiral-frontier normalisation
    $\mathcal{N}_{A}$~\cite{Beau2026aar}, and is not fixed in this note; what is fixed here is the assignment map
    $e_{0}, e_{+}, e_{-} \mapsto$ three ordered generation levels, prior to any identification with the physical
    $e, \mu, \tau$.

  \section{Status of results}
    \label{sec:status}

    \paragraph{Derived (this note).}
      The functorial Yukawa line $\LY = \wedge^{2}(\Sbundle)$ is the unique line of the carrier, invisible to the
      $\SU(2)_{L}$ adjoint sector and carrying the abelian weight (Proposition~\ref{prop:yukawa-line}).
      The generation-level assignment on $\Cgen$ is fixed by the spectral levels of
      $\Epi^{2}|_{\Cgen} = \diag(1, \tfrac12 + u, \tfrac12 - u)$ through the exit deficits
      (Proposition~\ref{prop:generation-level-assignment}).
      Both are structural given the carrier identifications of~\cite{Beau2026q14, Beau2026prs} and
      the algebraic normal form; no step invokes the conditional model of~\cite{Beau2026a34}.

    \paragraph{Open.}
      The projected Yukawa operator $Y_{\Pi}$ of~\eqref{eq:yukawa-operator}, its norm, the orientation (sign) of $u$, and
      the complex metaplectic mixing phase are not supplied here.
      Hence the absolute masses, the mass ratios beyond the dimensionless level structure, and the generation mixing
      remain open, in agreement with the open deliverables of~\cite{Beau2026fmnote}.

    \paragraph{Summary.}
      Front~3a is closed as a structural result: the Yukawa line and the generation ordering, not the mass values.
      The determinant line $\LY$ fixes the coupling line; the squared residue $\Epi^{2}|_{\Cgen}$ fixes the levels; the
      mass comes after.

  \section{Reproducibility}
    \label{sec:repro}

    All representation-theoretic and spectral identities are verified by exact symbolic computation, with no
    floating-point sampling, in \texttt{code/front3\_yukawa\_line.py} of this repository.
    The functorial dimensions $\dim \Sym^{2}(\mathbb{C}^{2}) = 3$ and $\dim \wedge^{2}(\mathbb{C}^{2}) = 1$, the
    $\SL(2,\mathbb{C})$ action on the line $\wedge^{2}(\mathbb{C}^{2})$ as $\det = 1$, and the abelian weight as the trace
    $y_{1} + y_{2}$ are checked symbolically (Proposition~\ref{prop:yukawa-line}).
    The normal form $\Epi^{2}|_{\Cgen} = \diag(1, \tfrac12 + u, \tfrac12 - u)$, its even-sector reduction to
    $\diag(1, \tfrac12, \tfrac12)$ at $u = 0$, the exit deficits, the ordering for $u > 0$, and the scale invariance of the
    level ratio are checked symbolically (Proposition~\ref{prop:generation-level-assignment}).

  \newpage

  \bibliographystyle{plain}
  \bibliography{cosmochrony-bibliography}

\end{document}
